%%%%%%%%%%%%%%%%%%%%%%%%%%%%%%%%%%%%%%%%%%%%%%%%%
%------ LaTeX-Template
%%%%%%%%%%%%%%%%%%%%%%%%%%%%%%%%%%%%%%%%%%%%%%%%%

%---- Header (mit Formateinstellugen) laden, Inputencoding prüfen ------

%%%%%%%%%%%%%%%%%%%%%%%%%%%%%%%%%%%%%%%%%%%%%%%%%
%---- LaTeX-Header
%%%%%%%%%%%%%%%%%%%%%%%%%%%%%%%%%%%%%%%%%%%%%%%%%

\documentclass[12pt,paper=A4,numbers=noenddot,bibliography=totoc,listof=totoc,DIV=11,BCOR=1mm]{scrreprt}
%  Bindekorrektur BCOR (verlorener Rand am linken Blattrand)! Wert haengt von der Art der Heftung ab!!
%  Satzspiegeleinstellung DIV von KOMA-Script / sccreprt.

\pagestyle{headings}

\usepackage[T1]{fontenc} % Font Encoding fuer europaeische Schriften mit Umlauten (Unterstuetzung der Worttrennung)
\usepackage{lmodern} % PostScript-Varianten der TeX Computer Modern-Schriften laden
\usepackage[english,ngerman]{babel} % Spracheinstellungen fuer Englisch und Neudeutsch laden

\usepackage{setspace}

\usepackage{graphicx} % Grafikeinbindung (fuer .JPG, .JPEG, .PNG und .PDF, falls pdflatex benutzt wird)
\usepackage[table]{xcolor} % ermoeglicht farbige Schrift und farbige Tabellenzeilen
\definecolor{black}{gray}{0} % Umdefinition der Farbe black, falls noetig (0=schwarz, 1=weiss)
\definecolor{dblue}{rgb}{0.1,0.2,0.6} % Dunkelblau, fuer Hyperlinks
\definecolor{lgray}{gray}{0.9} % Hellgrau, fuer Tabellen (0=schwarz, 1=weiss)

\usepackage{booktabs} % fuer schoene Tabellen

\usepackage[round,authoryear]{natbib} % Literaturverweise mit Name/Jahreszahl in runden Klammern
\bibpunct[:\,]{(}{)}{,}{a}{}{,~}  % Feinformatierung der Natbib-Zitierweise

\usepackage[hyphens]{url}
\usepackage[colorlinks=true,linkcolor=black,citecolor=dblue,urlcolor=dblue]{hyperref} 
\usepackage{hyperref}  
% die Pakete url und hyperref ermoeglichen anklickbare URLs im Quellenverzeichnis in definierter Farbe, 
% sie ermoeglichen den Zeilenumbruch bei langen URLs, und sie erzeugen Hyperlinks (Farbe s.o.) 
% zwischen Quellenverweis und Quellenverzeichnis sowie zwischen label und ref im PDF-Dokument

\usepackage[section]{placeins}
% Fonteinstellungen fuer Bildunterschriften: Unterschrift serifenlos, "Abbildung" fett (bfseries = bold face series)
\setkomafont{captionlabel}{\sffamily\bfseries}
\setkomafont{caption}{\sffamily}

%------------------------------------------------------------------------------------------------------------------
%------ Eigenstaendigkeitserklaerung im gerahmten Kasten (parbox in einer framebox) ------
%------------------------------------------------------------------------------------------------------------------

\newcommand{\eigen}{
\setlength{\fboxsep}{2ex}
\setlength{\fboxrule}{0.8pt} 
% Einstellungen fuer Rahmenabstand und Rahmendicke der Framebox
\begin{center}
	\fbox{
		\parbox{0.8\linewidth}{
		Ich versichere, die vorliegende Arbeit selbstst\"andig ohne fremde Hilfe verfasst 
		und keine anderen Quellen und Hilfsmittel als die angegebenen benutzt zu haben. 
		Die aus anderen Werken w\"ortlich entnommenen Stellen oder dem Sinn nach 
		entlehnten Passagen sind durch Quellenangaben eindeutig kenntlich gemacht.
		\par\bigskip\bigskip\bigskip\bigskip
		\hspace*{0.8cm}Ort, Datum \hfill \vorname~\nachname\hspace*{0.8cm}
		}
	}
\end{center}
}

%%%%%%%%%%%%%%%%%%%%%%%%%%%%%%%%%%%%%%%%%%%%%%%%%

%\usepackage[applemac]{inputenc} % Inputencoding f�r Mac
%\usepackage[latin1]{inputenc} % Inputencoding f�r PC/Win
\usepackage[utf8]{inputenc} % Inputencoding, universell
%\usepackage[utf8x]{inputenc} % Inputencoding, universell


%------------------------ Titelblatt-Layout laden ----------------------------------

%%%%%%%%%%%%%%%%%%%%%%%%%%%%%%%%%%%%%%%%%%%%%%%%%
%------ LaTeX-Titelblatt
%------ Deklarationen fuer die Titelseite
%%%%%%%%%%%%%%%%%%%%%%%%%%%%%%%%%%%%%%%%%%%%%%%%%

\title{\titel\\[2ex]
\LARGE Bachelor-Thesis\\
\large zur Erlangung des akademischen Grades B.Sc.\\[1.5ex]
\LARGE \vorname~\nachname\\[0.5ex] 
\large \matrikelnummer
}

\author{\unitlength1mm
\large\raisebox{-1ex}{\includegraphics[width=4em]{HAW_wuerfel}}\hspace{1ex}
\parbox[b]{11.2cm}{\sffamily\large%
Hochschule für Angewandte Wissenschaften Hamburg\\[-0.2ex]
Fakultät Design, Medien und Information\\[-0.2ex]
Department Medientechnik
}\\[6ex]
\sffamily\large Erstprüfer: \erstpruef\\[0.5ex]
\sffamily\large Zweitprüfer: \zweitpruef}

%%%%%%%%%%%%%%%%%%%%%%%%%%%%%%%%%%%%%%%%%%%%%%%%%

%---------------------------- Titeldefinitionen --------------------------------------

\newcommand{\vorname}{Sebastian}
\newcommand{\nachname}{Bohn}
\newcommand{\matrikelnummer}{2036605}

\newcommand{\titel}{Analyse und Evaluierung von plattformübergreifenden Spiel-Engines und Frameworks,anhand der Implementierung einer mobilen Beispielapplikation}

\newcommand{\erstpruef}{Prof. Dr. Edmund Weitz}
\newcommand{\zweitpruef}{Prof. Dr. Andreas Plaß}

\date{vorläufige Fassung vom \today}   % Vorab-Version 
%\date{\sffamily Hamburg, DD. MM. YYYY}  % Abgabedatum!

%----------------------------- ANFANG --------------------------------------

\begin{document}
\selectlanguage{ngerman}
% Titelseite erzeugen
\maketitle
% Inhaltsverzeichnis erzeugen          
\tableofcontents
% Seitenumbruch
\clearpage

%----------------------------INHALT---------------------------------------

%---------------------------ABSTRACT------------------------------------

\thispagestyle{empty}
\selectlanguage{english}
\section*{\centering\abstractname}

%TODO: English abstract

\selectlanguage{ngerman}
\section*{\centering\abstractname}

%TODO: Deutsches abstract

%----------------------------TEXT-----------------------------------------------

%TODO: Kapitel 1
\chapter{Einleitung}

\section{Motivation}
\section{Gliederung}

%TODO: Kapitel 2
\chapter{Mobile Systeme}

\section{Aktuelle Systeme auf dem mobilen Markt}
\subsection {iOS}
\subsection{Android}
\subsection{Windows Phone}
\subsection{Weitere Systeme}

\section{Bedarfsanalyse}
\subsection{Markt- und Useranteile der jeweiligen Systeme}
\subsection{Verfügbare Applikationen / Games der Stores}

%TODO: Kapitel 3
\chapter{Native Softwareentwicklung in den jeweiligen Systemen}
\section{Hardwarevorraussetzungen}
\section{Programmiersprachen}
\section{Entwicklungsumgebungen}

%TODO: Kapitel 4
\chapter{Cross-Plattform Entwicklung}
\section{Sinn und Gedanke von Cross-Plattform Entwicklung}
\section{Funktionsweise von Cross-Plattform Entwicklung}
\subsection{Technik}
\subsection{Geteilter Content}
\subsection{Übersetzung in die jeweiligen System}

%TODO: Kapitel 5
\chapter{Cross-Plattform Frameworks}
\section{Tools und Anbieter zur Entwicklung}
\section{Verweis auf Bachelorarbeit: „Plattformabhängige und –unabhängige Entwicklung mobiler Anwendungen am Beispiel von Geo-Wikipedia-App“}
\section{Gamespezifische Frameworks und Engines}
\subsection{Monogame}
\subsection{Cocos2D}
\subsection{Libgdx}
\subsection{Unity}
\subsection{Unreal Engine}
\subsection{Weitere Frameworks}
\section{Entwicklungsumgebungen}
\subsection{Unterstützte IDEs}
\subsection{Systembedingte Einschränkungen}

%TODO: Kapitel 6
\chapter{Gegenüberstellung der Frameworks}
%Features und Einschränkungen
\section{Zielplattformen}
\section{Skalierbarkeit der Menge der Plattformen} % ???
\section{Programmiersprachen}
\section{Unterstützung von 2D und 3D}
\section{Zugriff auf Hardware}%Accelerometer etc
\section{Free- und Pro- Versionen}
\section{Einfluss auf Einstellungen}
\section{Zusätzlich benötigte Software}
\section{Aktualität - Versionen - Community}
\section{Zukunftsaussichten}

%OPTIONAL!!!
%TODO: Kapitel 7
\chapter{Analyse der Marktanteile}
\section{Menge an Firmen und Entwicklern}
\section{Menge an Games}

%OPTIONAL!!!
%TODO: Kapitel 8
\chapter{Stores für mobile Spiele}
\section{Allgemeine Bedingungen für Entwickler}
\section{Maximale App-Größe}
\section{Anforderungen an den Quellcode}
\section{Kosten und Abgaben}

%TODO: Kapitel 9
\chapter{Kosten-Nutzen Vergleich}
%Vergleichstabelle zu den Variationen der Projektanforderungen im Zusammenhang mit den Möglichkeiten und Kosten der Frameworks

%TODO: Kapitel 10
\chapter{Grundgerüst und Aufbau eines Cross-Plattform Projekts}
\section{Geteilter Content}
\section{Plattformabhängiger Content}
\section{Grundaufbau bei Engines}
\section{Grundaufbau bei Frameworks}

%OPTIONAL!!!
%TODO: Kapitel 11
\chapter{Game-typische Design Patterns und Architekturen}
\section{Architekturen}
\section{Patterns}
\section{Verweis auf Bachelorarbeit: „Use of Design Patterns for mobile game Development“}

%OPTIONAL!!!
%TODO: Kapitel 12
\chapter{Nutzen von Architekturen}
\section{Pro}
\subsection{Skalierbarkeit}
\subsection{Lesbarkeit}
\subsection{Wiederverwertbarkeit}
\section{Contra}
\subsection{KISS - Keep it simple stupid}

%TODO: Kapitel 13
\chapter{Konzeption einer Applikation}
\section{Ideen}
\section{Anforderungen}
\section{User Stories}

%TODO: Kapitel 14
\chapter{Implementierung der Applikation}
\section{Verwendete Frameworks und Engines}
\section{Verwendete APIs und SDKs}
\section{Assets und deren Verwendung}

\chapter{Analyse messbarer Metriken}
%Kosten, Performance, Akkuverbrauch, App-Daten-Aufteilung, Größe des benötigten Speichers, Größe des geteilten Contents, Stabilität, RAM, Code-Zeilen/-Größe, Entwicklungszeit(bedingt zu berücksichtigen)

\chapter{Vergleich der Messprotokolle}

\chapter{Fazit}



%--------------------- VERZEICHNISSE -------------------------------------

\listoffigures % Abbildungsverzeichnis erzeugen
%\listoftables % Tabellenverzeichnis erzeugen

%------------------------------ LITERATURVERZEICHNIS----------------------
\begin{thebibliography}{}
\end{thebibliography}
%----------------------------- EIGENSTÄNDIGKEITSERKLÄRUNG-----------------
\clearpage\thispagestyle{empty}
\eigen  % im header definiert
%----------------------------- ENDE --------------------------------------
\end{document}